\subsection{Actuator Gating: {STPA} for Tool Calls}\label{sec:actuator}

The gate of Section~4 governs what the model \emph{reads}. Tool calls are a
different problem: they actuate. A softened refusal costs a turn; a deleted
directory costs data. Tool side effects are irreversible, and the hazards are not
injection-side---the model may emit a call while a conflict is unresolved, while
context usage sits at a preemptive threshold, or after a social-engineering
turn. We therefore extend the same causal-interposition idea to the interval
between a tool request and its execution, with a separate unsafe-control-action
namespace (\emph{T-UCA}) that neither modifies nor merges with the
injection-side UCA-01--08. The two gates are parallel: injection pre-flight,
then tool gate, then execution.

\paragraph{Scope, stated as an override.} Current project policy freezes
tool-loop work. This extension is an \emph{explicit override}, not an
evidence-driven unfreeze: the pre-declared trigger---real logs showing at least
ten tasks failed because external calls were unavailable, plus a demonstrated
inability to roll back tool errors---is \emph{not} met. We report it because it
is implemented and testable, but its existence is not evidence that the trigger
was satisfied. Three items remain frozen: multi-step tool planning,
transactional rollback or compensation, and writing tool outcomes into long-term
memory.

\paragraph{The T-UCA matrix.} Sixteen rules---contract version
\texttt{t-uca-1.1}, the paper-baseline snapshot (commit \texttt{d551079})---in six
groups. \emph{State coupling} (four): a mutating call
under variety fallback or unresolved conflict, in the critical phase, an
irreversible write with a saturated affective baseline, and a destructive call in
the shift phase. \emph{Data protection} (one): credentials in clear text.
\emph{Mutation safety} (four): destructive delete or overwrite without exact
confirmation, writes to frozen or protected paths, a mutating call after a failed
frozen-region check, and path escape. \emph{Network and execution} (two):
\mbox{non-allowlisted} egress and arbitrary code execution. \emph{Shape and registry}
(four): oversized or recursive payloads, schema or delimiter violations,
unregistered tools, undeclared capabilities. \emph{Auditability} (one):
incomplete audit context. Thirteen rules block, three degrade: a degrade is not
executed but requires explicit human confirmation through the host approval
seam, and falls back to block when no approval channel exists. No rule returns
\emph{allow}; allow is the default only for read-only calls that match nothing,
and unknown or undecidable calls default to block.

\paragraph{Why state matters.} Unlike the injection gate, which reads only the
prompt, the tool gate consumes host state: variety fallback, unresolved conflict,
the phase of the preemptive monitor, integral saturation, frozen-layer
verification, declared capabilities, and audit context. The same call is not the
same act in every state: deleting a file while a conflict is unresolved differs
from deleting it after clarification, and a write during the critical phase is
where an irreversible mistake is hardest to attribute.

\paragraph{Three attachment points.} We state which wirings are preventive.
(i)~In the engine's reference loop the gate sits after parsing and before
execution, so a blocked call is never executed. (ii)~In the host plugin,
\texttt{guardTool(spec, stateProvider)} wraps registered tools so that a call
routed through the wrapper is evaluated first: preventive, but only for tools
that route through it. (iii)~The host exposes \texttt{tool/call} and
\texttt{tool/result} session events but no pre-dispatch hook, so the plugin's
default path is \emph{detective}: it evaluates the call, records a
\texttt{tool\_gate} audit event (tool, argument digest, verdict, matched
identifiers in \texttt{ucas}, policy version, state snapshot), and injects a
strong clarification prompt on the next turn. Detection is not blocking, and we
label it as such.

\paragraph{Explainable permission transfer.} On a degrade the gate records a
permission request we call \emph{explanatory permission transfer}: the verdict
restated in the persona's voice as an explanation a non-expert can act on,
recomposed from the STPA gate's UCA rules, the persona baseline and the relation
state. The \emph{fact layer} fixes action, object, worst case, reversibility and
why now; these are mandatory and cannot be omitted or rewritten. The \emph{style
layer} maps the persona's pleasure--arousal--dominance (PAD) baseline and the
relation state onto tone, address and intensity by deterministic rules with zero
model calls. A composer injects both into the next turn's prompt; the model's
final wording is not verified. An \emph{executed} flag could mark the request as
held pending confirmation at a pre-dispatch point (engine loop,
\texttt{guardTool} routing), but the plugin wires only the detective path, where
the call may already have been dispatched; the request arrives as an explainable
account plus a request to continue or remediate, and no pre-dispatch hold is
claimed for arbitrary tools. Five T-UCA rules for delete, credentials, egress,
execution and path escape---together with any destructive or unclassified tool
class---form a red zone: styling is stripped entirely, persona
voice dropped, and the request becomes a serious, neutral warning, lest a joking
tone mislead. The intended benefit, lower human-in-the-loop cognitive load, is a
design argument, not a measurement: no user study, no blind approval rate, no
effect size. Verdicts are unchanged: at a pre-dispatch point a degrade is still
not executed, a degrade without an approval seam still falls back to block, and
no new approval seam is created.

\paragraph{Evidence.} The Python suite covers twenty constructed dangerous
scenarios (all blocked), eight benign scenarios (none blocked), state-gated and
fail-closed cases, and a twenty-case Python--JavaScript check in which verdict,
matched identifiers and downgrade field all agree. The JavaScript simulation runs
sixty-five checks; the Python suite reports 79 passing assertions. The reference
demo emits the JSON cited here
(\texttt{paper/\allowbreak{}ablation/\allowbreak{}results/\allowbreak{}actuator-\allowbreak{}demo-\allowbreak{}20261027.\allowbreak{}json}): twelve scripted
calls, four allowed and executed, eight blocked and none executed, with nine
distinct rules among the matched identifiers. It uses a mock executor with zero
model calls and zero real writes, so these are scenario counts showing
``judge before execute'', not an interception rate. The same run evaluates a
ten-item held-out set written blind---its author did not read the rule table, the
implementation, or the tests---frozen at
\texttt{paper/\allowbreak{}ablation/\allowbreak{}tool/\allowbreak{}heldout-\allowbreak{}tool-\allowbreak{}calls-\allowbreak{}v1.\allowbreak{}jsonl} (SHA-256
\texttt{\seqsplit{f68fc1b3f19fc37e43b9856096d29b9a55de7dee9bf664a9b9a9592e070576d4}}); ten
of ten calls are blocked with arguments only and ten of ten with context, and its
endpoints are placeholder domains, so a hit cannot be a domain coincidence. We
report verdicts, not a correctness rate: whether each call \emph{should} be
blocked awaits adjudication. The gateway is additive: injection-side rules are
untouched, and three suite failures that predate this work are
environment-related and unchanged.

\paragraph{What the numbers are not.} We do not convert these counts into a
defence rate: designers chose the constructed scenarios, and the held-out set is
ten cases from one blind author. The gate never executes payloads, so its attack
surface is parsing and its false-positive rate is uncalibrated; recovery is a
strong hint plus human confirmation, not silent denial.

\paragraph{Boundaries.} First, without a host pre-dispatch hook the plugin cannot
stop an arbitrary host tool; true host-level pre-blocking is an upstream request,
not a shipped capability. Second, if a tool executes outside the gateway---another
plugin, or a host path that bypasses the wrapper---the gate is inert: it is a
control point that can be adopted, not system-wide enforcement. Third, the
excluded items above remain frozen, so the paper makes no claim about closed-loop
tool behaviour.
